\begin{longtable}{p{2.7cm} p{2.2cm} p{1.5cm} p{6.5cm}}
\toprule
Campo & Tipo & Tamaño & Justificación \\
\midrule
\endhead

carne &
int &
4 B &
Llave primaria, requerida explícitamente por el caso de estudio \\

identificacion &
char[13] &
13 B &
Cubre el caso más ancho entre cédula física (9 dígitos) y pasaporte alfanumérico: 12 caracteres. \\

nombre &
char[21] &
21 B &
20 caracteres útiles, margen razonable para nombres compuestos. \\

apellidos &
char[41] &
41 B &
Un solo campo, como lo pide el enunciado: 40 caracteres útiles cubren primer y segundo apellido juntos. \\

fechaNacimiento &
struct Fecha &
12 B &
TAD propio (día, mes, año como \texttt{int}), reutilizable y con sus propias reglas de validez. \\

sexo &
enum (int) &
4 B &
Dominio cerrado paraa evitar inconsistencias de tamanos -> representado como entero en vez de texto libre. \\

carrera &
int &
4 B &
Referencia a un catálogo de carreras en vez de repetir el nombre completo en cada registro. \\

nivelAcademico &
enum (int) &
4 B &
Dominio cerrado (pregrado, grado, posgrado, etc). \\

creditosAprobados &
int &
4 B &
Cantidad entera, sin componente fraccionario. \\

promedio &
float &
4 B &
Requiere parte decimal. \\

estado &
int (bool a mano) &
4 B &
0 para inactivo, 1 para activo; sostiene el borrado lógico sin usar \texttt{bool} \\

correo &
char[41] &
41 B &
Margen suficiente para direcciones institucionales y personales comunes. \\

telefono &
char[9] &
9 B &
8 dígitos (formato costarricense) más el terminador nulo. \\

direccion &
char[81] &
81 B &
Texto libre de longitud variable, con un tope generoso para evitar problemas \\

\bottomrule
\end{longtable}
